\documentclass[journal]{IEEEtran}

\usepackage{cite}
\usepackage{amsmath,amssymb,amsfonts}
\usepackage{graphicx}
\usepackage{textcomp}
\usepackage{xcolor}
\usepackage{booktabs}
\usepackage{multirow}
\usepackage{url}
\usepackage{hyperref}
\usepackage{array}
\usepackage{tabularx}

\newcolumntype{L}[1]{>{\raggedright\arraybackslash}p{#1}}
\newcolumntype{C}[1]{>{\centering\arraybackslash}p{#1}}

\begin{document}

\title{Tactile Sensing for Dexterous Robot Hand Manipulation:\\A Comprehensive Survey}

\author{Terry~Taewoong~Um
\thanks{T. T. Um is with Cosmax (e-mail: terry.t.um@gmail.com).}}

\markboth{IEEE Transactions on Robotics, Vol.~XX, No.~X, 2026}%
{Um: Tactile Sensing for Dexterous Robot Hand Manipulation}

\maketitle

\input{sections/abstract}

\begin{IEEEkeywords}
Tactile sensing, dexterous manipulation, robot hand design, vision-language-action models, embodiment retargeting, sim-to-real transfer
\end{IEEEkeywords}

\input{sections/introduction}
\input{sections/tactile_sensors}
\input{sections/robot_hand_design}
\input{sections/data_representation}
\input{sections/learning_methods}
\input{sections/vla_models}
\input{sections/retargeting}
\input{sections/integrated_systems}
\input{sections/discussion}
\input{sections/conclusion}

\section*{Acknowledgment}
This work was seeded from joint research seminar materials between Seoul National University and Cosmax. The author would like to thank Junsoo Ha, Jihwan Song, Taejoon Park, and Incheol Jeong (PhD candidates) for preparing the seminars, and Professors Frank Chongwoo Park, Yong-Lae Park, and Kyujin Cho for their collaboration.

This project was built using the Harness skill by Minho Hwang, and the website structure was adapted from AI Trend Onboarding by Minho Hwang.

AI tools were used in the production of this work: Claude (Opus 4.6) for literature survey, content generation, and manuscript preparation, and the gemini-3-image-generation skill (Gemini 3 Pro Image) for figure generation. The author takes full responsibility for the accuracy and integrity of the content presented in this paper.

\bibliographystyle{IEEEtran}
\bibliography{references}

\end{document}
